\documentclass[11pt,a4paper]{article}
\usepackage[margin=2.2cm]{geometry}
\usepackage{amsmath,amssymb,booktabs,xcolor,enumitem,graphicx,float}
\usepackage[normalem]{ulem}
\usepackage[hidelinks]{hyperref}
\newcommand{\citet}[1]{Krnjaic et al.}
\setlength{\parskip}{0.6em}\setlength{\parindent}{0pt}
\definecolor{plc}{HTML}{888888}
\newcommand{\placeholder}[1]{\textcolor{plc}{\small Replaces the placeholder: #1}\par}
\newcommand{\tbd}{\textcolor{orange}{TBD}}
\sloppy
\title{Proposed change, Results and Discussion with 8 and 16 agents}
\date{17 November 2026, version 5 (both fleets, served share first, plain style)}
\begin{document}
\maketitle

\section*{How to read this}
Part~1 says what changes from version 4. Part~2 is the Results and Discussion of the Experiments
chapter (\texttt{5.Experiments.tex}) as they would read. All of Part~2 is new text. A grey line
above a paragraph names the placeholder it replaces. Numbers sit in the tables, and the prose says
what they mean. Part~3 has the questions. Nothing is changed in the thesis yet.

\section*{Part 1. What changes from version 4}

\begin{itemize}[leftmargin=1.2em]
\item \textbf{Fleet 32 added.} \emph{(Author, 17 November.)} A third column for 32 agents, filled once
  its runs are trained and tested. The centralised controller and the arm with messages are not
  trained at 32 agents. The prose for 32 agents follows the numbers.
\item \textbf{Both fleets.} 8 and 16 agents are reported side by side. At 8 agents the learned
  controllers keep up on most seeds. At 16 agents they fall behind. Together they show where the
  learned controllers stop working, which 16 agents alone cannot.
\item \textbf{Served share is the main measure for the learned controllers.} It reads simply, since
  1.0 means a fleet serves all the work that arrives. Service time is compared only for Token
  Passing. The reason is in ``Keeping up'' below. Even the learned seeds that pass the check leave
  part of the work behind, so their orders wait five to ten times longer than under Token Passing.
\item \textbf{Token Passing is the reference, not the winner.} It shows what a planner that keeps up
  gains from storage, and it sets the bar the learned controllers are measured against.
\item \textbf{Messages.} They help under fixed storage at both fleets, and not under the other two
  rules. The text says exactly that.
\item \textbf{Decision time at 32 and 64 agents.} The scaling probe is added to Costs. Token Passing
  stays the fastest, but it grows faster with the fleet than the decentralised controller does.
\item \textbf{Order time is the headline.} \emph{(Author, 17 November.)} The pick time stands beside it.
\item \textbf{Keep-up threshold stays at 0.8.} \emph{(Author, 17 November.)} Every result uses it.
\item \textbf{The load and what could have been done differently.} \emph{(Author, 17 November.)}
  The main loads stay, because every run used them. The Discussion says the load assumed every
  controller would keep up, why the learned controllers do not reach Token Passing, and what could
  have been done differently, such as longer training episodes.
\item \textbf{What was expected.} The new Discussion paragraph answers the predictions of the
  Expected outcome one by one.
\item \textbf{Plain style.} One finding per sentence, the numbers in four small tables.
\item \textbf{Still open.} Token Passing at 8 agents plays seeds 0 to 19, not the test seeds, until
  its replay runs (marked $^\dagger$ here only). The centralised controller under fixed storage at
  8 agents comes from Kaggle. The section-based arm's fixed policy under congestion-aware storage
  (16 agents) ends this afternoon. Fleet 32 is still training.
\end{itemize}

\section*{Part 2. The text}

\subsection*{Results}

\placeholder{One opening sentence that states the main finding once it is known.}
The results are for 8 and 16 agents, the fleets every controller could be trained on in the time
available. Demand storage makes Token Passing's orders 14\,\% faster at 8 agents and 20\,\% faster at
16. The learned controllers keep up on most seeds at 8 agents and fall behind at 16, so for them
the gain cannot be measured.

\paragraph{Keeping up.}
\placeholder{Which conditions pass the keep-up check. Report the backlog trace of one condition that
falls behind.}
Token Passing serves all the work at both fleets, under every rule (Table~\ref{tab:served}).

The learned controllers serve most of the work at 8 agents. They serve between 0.81 and 0.91 of it,
and each keeps up on 11 to 15 of the 20 seeds. The arm with messages under fixed storage does best,
at 0.96 and 18 seeds.

At 16 agents they fall behind. The section-based arm serves the most, between 0.54 and 0.69. The
centralised arm serves about half. The decentralised arm serves 0.12 under fixed storage and 0.66
under congestion-aware storage. Only 2 to 7 seeds keep up, all of them for the section-based or the
decentralised arm.

Passing the check does not mean serving everything. A seed passes when it serves at least 0.8 of
its work. The rest waits, and the wait grows through the run. So even on seeds that pass, a learned
controller's orders take 300 to 3200 steps, against 150 to 280 under Token Passing. Their service
times say more about the backlog than about storage, so the learned controllers are compared by how
much they serve.

\begin{table}[H]
  \centering
  \caption{Share of the arriving work that is served, mean over 20 test seeds. In brackets, the seeds
    that keep up. n.t., not trained at that fleet.}
  \label{tab:served}
  \footnotesize\setlength{\tabcolsep}{3pt}
  \begin{tabular}{l ccc ccc ccc}
    \toprule
    & \multicolumn{3}{c}{8 agents} & \multicolumn{3}{c}{16 agents} & \multicolumn{3}{c}{32 agents} \\
    \cmidrule(lr){2-4}\cmidrule(lr){5-7}\cmidrule(lr){8-10}
    controller & fix. & dem. & cng. & fix. & dem. & cng. & fix. & dem. & cng. \\
    \midrule
    Token Passing  & 0.98 (20)$^\dagger$ & 0.99 (20)$^\dagger$ & 0.99 (20)$^\dagger$ & 1.00 (20) & 1.00 (20) & 1.00 (20) & \tbd & \tbd & \tbd \\
    centralised    & \tbd      & 0.90 (15) & 0.81 (11) & 0.54 (0) & 0.52 (0) & 0.59 (0) & n.t. & n.t. & n.t. \\
    section-based  & 0.87 (14) & 0.81 (12) & 0.91 (14) & 0.63 (7) & 0.54 (0) & 0.69 (5) & \tbd & \tbd & \tbd \\
    decentralised  & 0.84 (12) & 0.87 (13) & 0.82 (11) & 0.12 (0) & 0.39 (0) & 0.66 (2) & \tbd & \tbd & \tbd \\
    with messages  & 0.96 (18) & 0.66 (9)  & 0.82 (13) & 0.37 (0) & 0.28 (0) & 0.34 (0) & n.t. & n.t. & n.t. \\
    \bottomrule
  \end{tabular}
\end{table}

\paragraph{Locks.}
Locks explain most of the decentralised arm's failure at 16 agents (Table~\ref{tab:locks}). A seed
locks when two agents hold the same task for 500 steps or more. Locks are rare at 8 agents. At 16
agents the decentralised arm locks on every seed under fixed and demand storage, with or without
messages. Token Passing never locks.

\begin{table}[H]
  \centering
  \caption{Test seeds that lock, out of 20.}
  \label{tab:locks}
  \footnotesize
  \begin{tabular}{l ccc ccc ccc}
    \toprule
    & \multicolumn{3}{c}{8 agents} & \multicolumn{3}{c}{16 agents} & \multicolumn{3}{c}{32 agents} \\
    \cmidrule(lr){2-4}\cmidrule(lr){5-7}\cmidrule(lr){8-10}
    controller & fix. & dem. & cng. & fix. & dem. & cng. & fix. & dem. & cng. \\
    \midrule
    Token Passing  & 0 & 0 & 0 & 0 & 0 & 0 & \tbd & \tbd & \tbd \\
    centralised    & \tbd & 0 & 0 & 5 & 3 & 2 & n.t. & n.t. & n.t. \\
    section-based  & 0 & 6 & 0 & 7 & 8 & 4 & \tbd & \tbd & \tbd \\
    decentralised  & 7 & 1 & 6 & 20 & 20 & 5 & \tbd & \tbd & \tbd \\
    with messages  & 0 & 10 & 5 & 19 & 20 & 18 & n.t. & n.t. & n.t. \\
    \bottomrule
  \end{tabular}
\end{table}

\paragraph{RQ1, storage adaptation.}
\placeholder{Gain in mean service time of $F_{\mathrm{dem}}$ and $F_{\mathrm{cng}}$ over
$F_{\mathrm{fix}}$ for each controller.}
Demand storage makes Token Passing's orders faster at both fleets (Table~\ref{tab:gain}). The gain is
14\,\% at 8 agents and 20\,\% at 16, and it shows on nearly every seed. A single pick gains less, 15\,\% at 16 agents [11.3, 18.3]. Fleet 8's pick gain follows once its test-seed rows are home. The time saved is travel and
the wait for an agent. Popular stock sits closer to the workstations, so a trip is shorter and the
agent is free sooner. Time spent blocked stays near 2 steps per pick.

For the learned controllers demand storage has no clear effect on how much they serve. It raises
the share for some and lowers it for others (Table~\ref{tab:served}). Each rule has its own trained
policy, so a change here mixes the rule with the training run.

\begin{table}[H]
  \centering
  \caption{Token Passing's mean order service time under fixed storage, and the gain of each rule
    over it, with 95\,\% intervals over 20 seeds paired by seed.}
  \label{tab:gain}
  \small
  \begin{tabular}{l c cc cc}
    \toprule
    fleet & $\bar\zeta_T$ under $F_{\mathrm{fix}}$ & gain of $F_{\mathrm{dem}}$ & seeds faster
      & gain of $F_{\mathrm{cng}}$ & seeds faster \\
    \midrule
    8 agents$^\dagger$  & 177.5 & 14.1\,\% [8.3, 19.9]  & 18 & 12.3\,\% [6.7, 18.0]  & 17 \\
    16 agents           & 278.9 & 20.0\,\% [16.0, 24.1] & 20 & 19.9\,\% [15.0, 24.9] & 20 \\
    32 agents           & \tbd & \tbd & \tbd & \tbd & \tbd \\
    \bottomrule
  \end{tabular}
\end{table}

\paragraph{RQ1, the controller.}
\placeholder{Compare the gains of Table~\ref{tab:gain} across controllers.}
The expected order, with the centralised controller first, does not hold. At 8 agents the three
learned controllers serve about the same. At 16 agents the section-based controller serves the
most and the decentralised controller the least, except under congestion-aware storage, where the
two are close. Deciding for the whole fleet at once did not pay.

\paragraph{RQ2, congestion-aware storage.}
\placeholder{Change in crowding $\bar\delta_T$ and blocked time $W_T$ from demand to congestion-aware
storage.}
Congestion-aware storage gives Token Passing no more than demand storage does
(Table~\ref{tab:gain}). Token Passing plans each path ahead, and an agent waits its turn as part of
the plan. Its agents are blocked for under 2 steps per pick, so the rule has little to avoid.

The learned controllers are blocked far more, 13 to 26 steps per pick at 16 agents. There, every
learned controller serves more under congestion-aware storage than under demand storage. The gain
comes with the training, not with the rule alone. The policy trained under fixed storage, played
under congestion-aware storage, serves no more than under fixed storage (Appendix~L). At 8 agents
the effect is mixed.

\paragraph{RQ4, communication.}
\placeholder{Only if reached.}
Messages help the decentralised controller under fixed storage, at both fleets. At 8 agents its
served share rises from 0.84 to 0.96, the highest of any learned run, and at 16 agents from 0.12 to
0.37. Under demand storage messages lower the share at both fleets. Under congestion-aware storage
they lower it at 16 agents and leave it unchanged at 8. At 16 agents the arm with messages still
locks on 18 to 20 seeds, so messages do not let two agents back out of a lock.

\paragraph{Costs.}
\placeholder{Mean service time of put-aways, runtime per controller.}
Demand storage costs put-away time. Under Token Passing at 16 agents a put-away takes 37 steps under
fixed storage and 43 under demand storage, because new stock goes further to reach the best shelves.

Decision time grows with how much of the warehouse a controller reads (Table~\ref{tab:time}).
Token Passing is the fastest at every fleet size. It also grows the fastest, about tenfold from 16
to 64 agents, while the decentralised controller grows less than threefold. The 32 and 64 agent
times come from policies trained at 16 agents and played for 500 steps, so they measure time only,
not how well the controllers do.

\begin{table}[H]
  \centering
  \caption{Mean time per decision step in milliseconds. Learned controllers on the GPU, Token
    Passing on the CPU.}
  \label{tab:time}
  \small
  \begin{tabular}{l ccc}
    \toprule
    controller & 16 agents & 32 agents & 64 agents \\
    \midrule
    Token Passing  & 2   & 6$^\ast$   & 19$^\ast$ \\
    centralised    & 110 & 266 & 792 \\
    section-based  & 73  & 106 & 235 \\
    decentralised  & 13  & 15  & 32 \\
    \bottomrule
  \end{tabular}

  \smallskip
  {\footnotesize $^\ast$ At the lowest load tested, 0.4 and 0.8 tasks per step, above the main load.}
\end{table}

\subsection*{Discussion}

\placeholder{Open with the answer to the main research question in one or two sentences.}
Storage that follows demand makes orders 14 to 20\,\% faster, when the controller keeps up. Only
Token Passing keeps up on every seed, so only for Token Passing can the gain be measured.

\paragraph{What was expected.}
\placeholder{The TODO added on 17 November.}
Some predictions of the Expected outcome held and some did not. Demand storage did shorten orders,
through shorter trips, as predicted. The learned controllers did not keep up at 16 agents and did
not come close to Token Passing. The centralised controller did not gain the most. It never kept up
at 16 agents, and at 8 agents it served about as much as the others. Closing the loop changed little for
Token Passing, as predicted, because it is rarely blocked. Messages helped only under fixed
storage, not where routes cross in general. RQ3 waits for the larger fleets.

\paragraph{The load.}
The main load assumed that every controller could keep up. It was set from what Token Passing can
serve, at about 0.6 of its capacity, before any learned controller was trained
(Appendix~E). Under the served check of the final runs, Token Passing's capacity is
one tested step higher, so the main load sits at about 0.5 of it. Every run was trained and tested
at this load, so it stays. The assumption held for Token Passing and nearly held for the learned
controllers at 8 agents. At 16 agents it did not hold.

\paragraph{Why adaptation helps.}
Demand storage helps because it puts popular stock near the workstations. A pick travels less and
its agent is free sooner. Blocked time hardly moves, so there is no sign that adaptation crowds the
aisles.

\paragraph{Why the learned controllers fall behind.}
The learned controllers fall behind mostly through locks and through runs that were too short.
Two agents that each wait for the other hold one task for the rest of the run, and the agents behind
them queue. Krnjaic et al.\ saw the same with agents that share one policy, and a manager that spread
the goals over the warehouse avoided it. The section-based controller does much the same and locks
less. The policies also trained and were chosen on runs of 2000 steps, while a test lasts 10000. A
policy that falls behind slowly looks fine in the short run.

\paragraph{What could have been done differently.}
Four changes could have brought the learned controllers closer to Token Passing.
\begin{itemize}[leftmargin=1.2em]
  \item \emph{Longer training episodes.} The policies trained on 2000 steps and were tested on
    10000. A backlog that grows slowly hardly shows in 2000 steps, so training never punished it.
    Training on longer episodes, or choosing the checkpoint on a long validation run, targets the
    failure seen in the test.
  \item \emph{More training.} Each run had 84 iterations and one seed. More iterations, or several
    training seeds per controller and rule, would show whether a run is simply unlucky.
  \item \emph{A way out of locks.} Most of the decentralised arm's loss is locks. A shared reward,
    or a rule that makes one of two locked agents give way, would let them back out.
  \item \emph{Starting from Token Passing.} Training the policies first to copy Token Passing's moves,
    as PRIMAL does with a planner, would start them from a fleet that already keeps up.
\end{itemize}
Longer test runs would not have helped. A learned run that falls behind falls further behind the
longer it runs, so a longer test would only lower its served share.

\paragraph{The price of avoiding congestion.}
Congestion-aware storage pays only where agents get blocked. Token Passing is rarely blocked, so the
rule has nothing to win there. The learned controllers are blocked often, and the policies trained
under the rule serve more at 16 agents. Whether the rule pays for a learned controller that keeps up
stays open.

\paragraph{Relation to earlier work.}
Learned controllers beat hand-made rules in Krnjaic et al.\ and lose to Token Passing here. The two
results do not clash. Their agents choose only where to go next, and a path search moves them there.
Here the learned controllers choose every move and are compared with a planner whose paths never
meet. Learning which task to take can pay, while learning to move 16 agents without collisions is
where the learned controllers fall behind. Lifelong MAPD, the GCN routing of van Knippenberg et al.\
and crowd-aware routing, still to write.

\paragraph{Threats to validity.}
Three items added.
\begin{itemize}[leftmargin=1.2em]
  \item The checkpoint is chosen on runs of 2000 steps and tested on 10000. Every learned policy did
    better in validation than in the test.
  \item Each controller and rule has one training run. A rule whose policy does better may only have
    had a better run. Playing one policy under every rule separates the two, and there the rule alone
    made no difference.
  \item The decision times at 32 and 64 agents come from policies trained at 16 agents.
\end{itemize}

\subsection*{Appendix L}
Tables~\ref{tab:served} and \ref{tab:locks} also go in Appendix~L, ``The test runs in full'', with the
lock threshold of 1250 steps beside 500.

\section*{Part 3. Questions}

Settled on 17 November. The keep-up threshold stays at 0.8. Token Passing at 8 agents is replayed
on the test seeds (running). Fleet 32 gets its own column. The order time is the main result, with
the pick time beside it (at 16 agents, orders 20.0\,\% faster and picks 14.8\,\% faster).

No questions are open.

\end{document}
